\section{Introduction}

Static analysis tools --- ruff for style, mypy for types --- catch fewer than one in three
of the code generation failures that stump GPT-4o-mini on EvalPlus. On 34 HumanEval+
failures, ruff and mypy fire on 32.4\% of cases; on 100 MBPP+ failures, the fire rate
drops to 16.0\%. The errors that remain are not syntax bugs. They are semantic
divergences: the model's algorithm is logically wrong, and asking it to ``please try
again'' without explaining \textit{what} went wrong provides almost no repair signal.

This observation is not a marginal footnote. It exposes a structural mismatch between
how most LLM repair pipelines are designed and what the dominant failure mode of modern
code generation models actually is. Current systems either provide no feedback
(baseline re-inference), raw error tracebacks that describe syntactic symptoms rather
than semantic causes, or broad self-critique loops that lack the precision to identify
algorithmic errors. The result is that the most common LLM code failures --- the ones
that survive round-0 evaluation --- are also the ones most poorly served by existing
repair oracles.

The deeper problem is one of information: targeted algorithmic repair requires three
things simultaneously. The model needs to know (1) what algorithm it was \textit{supposed}
to implement --- the formal intent captured in the problem docstring; (2) a concrete
behavioral gap --- a failing test's input and expected output, which acts as a
counterexample that specifies exactly where the model's behavior diverges from the
specification; and (3) a deviation signal --- the model's actual incorrect output,
enabling it to detect where its reasoning went wrong. No single component is
sufficient. A docstring without evidence of failure provides no diagnostic. A failing
test input without the expected output provides no target. The actual output without
the docstring leaves the model without a goal. The triple is a structured, jointly
motivated oracle for semantic repair.

This structured triple --- (docstring formal intent, failing test I/O pair, actual model
output) --- is what we call \textit{specification-aligned repair}. The insight motivating the
design is convergent: docstring re-anchoring is the primary driver of specification
grounding improvement \citep{Haeri2026}; contrastive I/O pairs outperform
single failure messages for automated program repair \citep{Kong2024}; and content,
not mere re-exposure, drives LLM repair improvement \citep{Iscan2026}. Taken together,
these three mechanisms are mutually reinforcing: the triple provides intent,
behavioral gap, and diagnostic simultaneously.

Despite this convergent motivation, no prior work has directly tested specification-aligned
repair against blind reprompting on EvalPlus's own semantic failure set using a
controlled 3-condition paired design. This gap matters because EvalPlus's augmented
test suite --- 80$\times$ more test cases than base HumanEval --- makes it a significantly stricter
evaluation environment than the benchmarks used in prior repair work, and because the
failure set from GPT-4o-mini round-0 represents a natural, unselected population of
semantic errors rather than a curated bug corpus.

This paper is a complete contribution in the infrastructure-and-preregistration
tradition of open science: we establish, verify, and publish the data infrastructure
and experimental design prior to executing any mechanism experiments, following the
principle that pre-registered designs provide stronger scientific guarantees than
post-hoc analysis \citep{Nosek2018}. The infrastructure and design are
independently valuable scholarly contributions to reproducibility and transparency in
LLM repair research.

We address this gap through the following contributions:

\textbf{Data infrastructure verification (this paper).} We establish and verify the data
infrastructure for reproducible specification-aligned repair experiments on EvalPlus.
The 134-problem GPT-4o-mini failure set from h-e1 Run 2 (34 HumanEval+ + 100 MBPP+)
is fully recoverable from archive: all 134 failure IDs are confirmed, all 134 stored
incorrect solutions are present in \texttt{solutions\_cache.jsonl}, the EvalPlus API is
accessible for all 134 task IDs, and deterministic test selection via \texttt{plus\_input[0]}
is confirmed (780 augmented test cases for a sample task). Six tasks with empty
\texttt{plus\_fail\_tests} fields define a 128-task working set with statistically adequate
power for the mechanism experiments.

\textbf{Pre-registered 3-condition comparative design.} We pre-register a McNemar-based
3-condition experiment --- Condition A (round-0 baseline, no repair), Condition B
(blind reprompting, ``please try again''), Condition C (specification-aligned repair,
structured triple) --- on the 128-task working set. The pre-registered predictions are:
(P1) Condition C achieves significantly higher round-1 pass@1 than Condition B
(McNemar one-tailed p $<$ 0.05); (P2) Condition C achieves significantly higher pass@1
than Condition A with fix rate $\geq$ 15\%; (P3) HumanEval+ fix rate exceeds MBPP+ fix
rate under Condition C (directional, exploratory).

\textbf{Methodological contribution.} We verify and document the data infrastructure for
this specific failure set, providing a reusable template for pre-experiment
reproducibility checks in LLM repair research --- enabling repair experiments
without new baseline API calls.

The remainder of the paper is organized as follows. Section~2 discusses related work
in LLM self-repair, specification grounding, and statistical repair evaluation design.
Section~3 describes the specification-aligned repair methodology and experimental
infrastructure. Section~4 details the experimental setup. Section~5 presents the
existence verification results and the pre-registered comparative design. Section~6
discusses implications and limitations. Section~7 concludes.
